\chapter{Statistical Analysis}
\label{chap:statistical_analysis}
\section{Analytical Aims and Data Provenance}
This chapter analysis the outcomes produced by the architectural design in \Cref{chap:arch_design} and the experimental design in \Cref{chap:experimental_design}. The primary aim is to answer the research questions outlined in \Cref{sec:research_questions}. The analysis will focus on the following key questions:
\begin{enumerate}[label=\textbf{RQ\arabic*:}, leftmargin=*]
    \item How do \emph{overlap ratio} and \emph{background noise} affect ASR accuracy under controlled synthetic conditions?
    \item How do publicly available ASR models compare across different overlap and noise operating points?
    \item To what extent do trends observed on controlled \emph{synthetic mixtures} transfer to real conversational recordings with natural overlap?
    \item Does the proposed double-stream serialisation metric provide a more suitable evaluation than ordinary WER when a two-speaker reference is compared with a single-stream hypothesis?
\end{enumerate}
The analysis uses the top-level result files generated by the evaluation pipelines described in \Cref{chap:arch_design}, which contain the WER, DSS-WER, ORC-WER and other relevant metrics for each ASR model across all experimental conditions. 

Two analysis scopes are kept separate:
\begin{enumerate}
    \item A 700 clips for each of the 4 ASR models. They are a cross-model subset of the 4800 clips generated for the synthetic evaluation, defined as
    \begin{equation}
        \begin{split}
        \mathcal{S}_{\mathrm{model}} ={}&
        \{\mathrm{SNR}=7.4,\ \mathrm{noise}=T,\ \mathrm{OVR}\in\{0.00,0.14,0.20,0.40\}\}\\
        &\cup
        \{\mathrm{OVR}=0.14,\ \mathrm{noise}=T,\ \mathrm{SNR}\in\{\mathrm{clean},7.4,0,-5\}\}.
        \end{split}
    \end{equation}
    Each combination of factors contains 100 clips with a total of 700 clips per model as shown in \Cref{fig:stat_cross_model_scope_support}. 
    \item the full set of 4800 clips for the \texttt{wav2vec-2.0} model, which is the only model evaluated across all conditions.   
\end{enumerate}

The first scope allows direct comparison of the 4 ASR models across the same set of clips. The reason for not using the full 4800 clips for all models is due to runtime and computational constraints. This is suitable for fair model ranking but should not be used to make claims about effects of SNR-overlap combinations for every models. The second scope allows a more comprehensive analysis of the effects of overlap and noise on the ASR performance of \texttt{wav2vec-2.0}, which has the best balance of performance and runtime among the models evaluated.

\begin{figure}[htbp]
    \centering
    \includefigure[width=0.68\linewidth]{figures/stat_cross_model_scope_support.png}
    \caption{Unique mixture support in the cross-model comparison scope. Each included SNR--OVR condition contains 100 mixtures; grey cells are outside the planned matched subset.}
    \label{fig:stat_cross_model_scope_support}
\end{figure}

\section{Variable Selection}
\subsection{Dependent Variables}
The primary dependent variable is the clip-level Word Error Rate (WER). The reason for choosing WER is that it is the de-facto metric for ASR evaluation \cite{sasindran2023hevalnewhybridevaluation} and it is available for all four models including segmented and non-segmented hypothesis.

ORC-WER and DSS-WER are also available where ORC-WER is only available for the segmented hypothesis (from \texttt{faster-whisper} and \texttt{whisperX}) and DSS-WER is available for all models. They act as secondary metrics to provide additional insights into the performance. They are valuable as they are designed to handle different situations. ORC-WER is designed to represent the speaker-agnostic lexical accuracy of the segmented hypothesis while DSS-WER is designed to evaluate the specific case of comparing a two-stream reference with a single-stream hypothesis, which is the case for this evaluation. However, since it is proposed in this dissertation, it is important to evaluate its behaviour and suitability before applying it as the primary metric. Therefore, WER is the most suitable and thus chosen as the primary metric for model comparison, while DSS-WER will be tested for its suitability and insights as a secondary metric.

\subsection{Explanatory Variables}
The main explanatory variables are signal-to-noise ratio (SNR), overlap ratio (OVR), noise type and ASR model. SNR and OVR represent the two main factors, noise and overlap, of interest in ASR evaluation. Overlap is the central phenonmenon being investigated in this dissertation. Noise type is included as it might interact with the effects of SNR and OVR. ASR model is included as it is important understand how different models perform under the same conditions, and whether the effects of SNR and OVR affect different models differently.

The speaker count is held constant at 2 for all clips. The makes OVR more interpretable as a controlled measure of two-speaker overlap rather than a more complex interaction between overlap and speaker count. The clip duration is also held at 60 seconds, which support comparability across clips and models, containing enough speech content while keeping the inference time manageable.

\section{Statistical Methods}
All reported results use the clip as the unit of analysis. The design is a factorial design with repeated measures, where same base audio is used across different SNR and noise type conditions, cross-model analysis uses the same clip IDs system across models. For this reason, the analysis will focus on matched condition summaries and comparisons.

WER is non-negative and often skewed, so the mean is reported as the primary summary statistic, but the median and standard deviation are also reported to provide a more complete picture.

\section{Overview of the Evaluation Results}
\subsection{Scope 1: Cross-Model Comparison}
\label{sec:cross_model_comparison}

\begin{table}[htbp]
\label{tab:cross_model_metric_availability}
\caption{Cross-model availability of WER, ORC-WER and DSS-WER.}
\centering
\begin{tabular}{lccc}
\toprule
Model & WER & ORC-WER & DSS-WER \\
\midrule
faster-whisper & yes & yes & yes \\
WhisperX & yes & yes & yes \\
wav2vec-2.0 & yes & no & yes \\
Parakeet & yes & no & yes \\
\bottomrule
\end{tabular}
\end{table}

\begin{table}[ht]
\centering
\caption{WER-first model comparison on the matched 700-clip subset.}
\label{tab:cross_model_wer_first}
\begin{tabular}{lrrrrr}
\toprule
\textbf{Model} & \textbf{Mean WER} & \textbf{Median WER} & \textbf{SD}  & \textbf{Mean ORC} & \textbf{Mean DSS}\\
\midrule
faster-whisper  & 0.159 & 0.123 & 0.128 & 0.604 & 0.156 \\
WhisperX  & 0.181 & 0.139 & 0.150 & 0.584 & 0.179\\
wav2vec-2.0 & 0.228 & 0.205 & 0.133 & - & 0.211\\
Parakeet  & 0.503 & 0.525 & 0.208  & - & 0.515\\
\bottomrule
\end{tabular}
\end{table}



\Cref{tab:cross_model_metric_availability} shows the availability of each metric across the different models. ORC-WER is not available for\texttt{wav2vec-2.0} and \texttt{Parakeet} as OCR-WER is designed for segmented hypothesis while wav2vec-2.0 and Parakeet output one continuous sequence. \Cref{tab:cross_model_wer_first} shows the mean, median and standard deviation of WER for each of the models. The results show that \texttt{faster-whisper} has the lowest mean WER, while \texttt{Parakeet} has the highest mean WER. 


\begin{figure}[htbp]
    \centering
    \includefigure[width=0.95\linewidth]{figures/stat_cross_model_wer_box.png}
    \caption{Violin plot of WER across different ASR models.}
    \label{fig:stat_cross_model_wer_box}
\end{figure}

\subsection{Inference Time}
The timing benchmark uses 100 randomly selected clips from the same synthetic corpus. Time per
clip is computed as total processing time divided by the number of clips. RTFx is computed as
the total audio duration divided by total processing time; a larger value therefore indicates
faster-than-real-time processing. In \Cref{fig:stat_time_per_clip_bar} and \Cref{tab:timing_efficiency}, \texttt{wav2vec-2.0} is the fastest model with a mean inference time of 0.754 seconds per clip, while \texttt{faster-whisper} is the slowest with a mean inference time of 27.736 seconds per clip. The RTFx values indicate the length of audio that can be processed in one second of real time, with \texttt{wav2vec-2.0} achieving the highest RTFx of 88.951. As shown in \Cref{fig:stat_accuracy_latency_tradeoff_panel}, the combination of relatively low WER and high RTFx makes \texttt{wav2vec-2.0} the candidate for the second scope of analysis, which focuses on the effects of SNR and OVR on a single model.

\begin{table}[htbp]
\centering
\caption{Inference time and real-time factor on the 100-clip timing benchmark.}
\label{tab:timing_efficiency}
\begin{tabular}{lrrr}
\toprule
\textbf{Model} & \textbf{Clips} & \textbf{Time per clip (s)} & \textbf{RTFx} \\
\midrule
wav2vec-2.0 & 100 & 0.754 & 88.951 \\
Parakeet & 100 & 1.124 & 59.663 \\
WhisperX & 100 & 7.276 & 9.220 \\
faster-whisper & 100 & 27.736 & 2.419 \\
\bottomrule
\end{tabular}
\end{table}

\begin{figure}[htbp]
    \centering
    \begin{subfigure}[t]{0.49\linewidth}
        \centering
        \includefigure[width=\linewidth]{figures/stat_time_per_clip_bar.png}
        \caption{Inference time per clip.}
        \label{fig:stat_time_per_clip_bar}
    \end{subfigure}
    \hfill
    \begin{subfigure}[t]{0.49\linewidth}
        \centering
        \includefigure[width=\linewidth]{figures/stat_accuracy_latency_tradeoff.png}
        \caption{WER--latency trade-off.}
        \label{fig:stat_accuracy_latency_tradeoff_panel}
    \end{subfigure}
    \caption{Inference efficiency across ASR models. The left panel reports mean inference time per approximately 60-second clip, while the right panel compares timing with matched-subset WER.}
\end{figure}

\section{Scope 2: Full Condition Analysis of wav2vec-2.0}

\begin{figure}[htbp]
    \centering
    \includefigure[width=0.82\linewidth]{figures/stat_wav2vec2_snr_overlap_heatmap.png}
    \caption{WER heatmap for wav2vec-2.0 across different SNR and overlap conditions.}
    \label{fig:stat_wav2vec2_snr_overlap_heatmap}
\end{figure}
\Cref{fig:stat_wav2vec2_snr_overlap_heatmap} shows the WER heatmap for \texttt{wav2vec-2.0} across different SNR and overlap conditions. There are in total 48 conditions 
\begin{equation}
    \mathcal{S}_{\mathrm{full}} = \{\mathrm{OVR}\in\{0.00,0.14,0.20,0.40\},\ \mathrm{noise}\in\{D, P, T\},\ \mathrm{SNR}\in\{\mathrm{clean},7.4,0,-5\}\},
\end{equation}
where each condition contains 100 clips. The heatmap shows a clear tread of increasing WER with higher overlap and lower SNR (noisier conditions).

This section went through the coverage and reliability of the evaluation outputs. The following sections analyse the results by research question, starting with the effects of overlap ratio and signal-to-noise ratio.

\section{RQ1: Effects of Overlap Ratio and Signal-to-Noise Ratio}

\begin{figure}[htbp]
    \centering
    \begin{subfigure}[t]{0.49\linewidth}
        \centering
        \includefigure[width=\linewidth]{figures/stat_wav2vec2_wer_vs_ovr.png}
        \caption{WER against overlap ratio, stratified by SNR.}
        \label{fig:stat_wav2vec2_wer_vs_ovr}
    \end{subfigure}
    \hfill
    \begin{subfigure}[t]{0.49\linewidth}
        \centering
        \includefigure[width=\linewidth]{figures/stat_wav2vec2_wer_vs_snr.png}
        \caption{WER against SNR condition, stratified by OVR.}
        \label{fig:stat_wav2vec2_wer_vs_snr}
    \end{subfigure}
    \caption{Mean wav2vec-2.0 WER trends across the full synthetic sweep. Each panel includes condition-specific lines and an overall pooled mean.}
    \label{fig:stat_wav2vec2_wer_trends}
\end{figure}



\begin{figure}[htbp]
    \centering
    \begin{subfigure}[t]{0.49\linewidth}
        \centering
        \includefigure[width=\linewidth]{figures/stat_wav2vec2_delta_vs_ovr.png}
        \caption{Delta WER against OVR, stratified by SNR.}
        \label{fig:stat_wav2vec2_delta_vs_ovr_panel}
    \end{subfigure}
    \hfill
    \begin{subfigure}[t]{0.49\linewidth}
        \centering
        \includefigure[width=\linewidth]{figures/stat_wav2vec2_delta_vs_snr.png}
        \caption{Delta WER against SNR, stratified by OVR.}
        \label{fig:stat_wav2vec2_delta_vs_snr_panel}
    \end{subfigure}
    \caption{Delta WER for wav2vec-2.0. Delta (a) is computed relative to the clean condition at the same overlap ratio  while Delta (b) is computed relative to the zero overlap condition at the same SNR.}
    \label{fig:stat_wav2vec2_delta_vs_ovr}
\end{figure}
\subsection{Effect of Overlap Ratio}
\Cref{fig:stat_wav2vec2_wer_vs_ovr} shows the general trends of WER against overlap ratio. The results shows that there is a significant increase in WER as the overlap ratio increases. 

\Cref{fig:stat_wav2vec2_delta_vs_snr_panel} shows the delta WER $\Delta_\text{WER, 0\% OVR}$ against SNR, where
\begin{equation}
    \Delta_\text{WER, 0\% OVR} = \text{WER}_{\mathrm{SNR}, 0\%} - \text{WER}_{\mathrm{SNR}, \mathrm{OVR}}.
\end{equation}
This isolate the effect of overlap as SNR changes. The results show that the $\Delta_\text{WER, 0\% OVR}$ increases as the SNR decreases, which suggests that the effect of overlap is more significant when SNR is low (noisier conditions). The $\Delta_\text{WER, 0\% OVR}$ diverges across different overlap ratio conditions as the SNR decreases, which matches the observation from \Cref{fig:stat_wav2vec2_delta_vs_ovr_panel} that the effect of overlap becomes more significant as the SNR decreases.

\subsection{Effect of Signal-to-Noise Ratio}
\Cref{fig:stat_wav2vec2_wer_vs_snr} shows the general trends of WER against SNR. The results shows that there is a significant increase in WER as the SNR decreases (gets noisier).
\Cref{fig:stat_wav2vec2_delta_vs_ovr_panel} shows the delta WER $\Delta_\text{WER, clean SNR}$ against overlap ratio, where
\begin{equation}
    \Delta_\text{WER, clean SNR} = \text{WER}_{\mathrm{SNR}, \mathrm{OVR}} - \text{WER}_{\mathrm{clean}, \mathrm{OVR}}.
\end{equation}
This isolate the effect of noise as overlap ratio changes. The results show that the $\Delta_\text{WER, clean SNR}$ decreases as the overlap ratio increases, which suggests that the effect of noise is more significant when overlap ratio is low, and decreases as the overlap ratio increases. The $\Delta_\text{WER, clean SNR}$ starts to converge across different SNR conditions as the overlap ratio increases, which suggests that the effect of noise starts to become dominated by the effect of overlap as the overlap ratio increases.

\Cref{fig:stat_wav2vec2_delta_vs_snr_panel} shows the delta WER $\Delta_\text{WER, 0\% OVR}$ against SNR, where
\begin{equation}
    \Delta_\text{WER, 0\% OVR} = \text{WER}_{\mathrm{SNR}, 0\%} - \text{WER}_{\mathrm{SNR}, \mathrm{OVR}}.
\end{equation}
This isolate the effect of overlap as SNR changes. The results show that the $\Delta_\text{WER, 0\% OVR}$ increases as the SNR decreases, which suggests that the effect of overlap is more significant when SNR is low (noisier conditions). The $\Delta_\text{WER, 0\% OVR}$ diverges across different overlap ratio conditions as the SNR decreases, which matches the observation from \Cref{fig:stat_wav2vec2_delta_vs_ovr_panel} that the effect of overlap becomes more significant as the SNR decreases.

To summarise the trend as a continuous estimating equation, SNR is converted to an amplitude noise ratio
\begin{equation}
    z(\mathrm{SNR}) =
    \begin{cases}
        0, & \mathrm{clean}\;(\mathrm{SNR}=\infty),\\
        10^{-\mathrm{SNR}_{\mathrm{dB}}/20}, & \text{otherwise},
    \end{cases}
\end{equation} 
changing SNR from logarithmic to linear scale, making SNR a continuous variable. Fitting a least-squares surface to the full \texttt{wav2vec-2.0} scope gives
\begin{equation}
    \widehat{\mathrm{WER}} =
    0.029 + 1.001\,\mathrm{OVR} + 0.274\,z - 0.328\,(\mathrm{OVR}\times z).
\end{equation}
The interaction term is retained because the noise penalty is not constant across overlap ratios. This equation should be interpreted as an estimate of the mean WER averaged over noise type: it gives a strong fit to the 16 SNR--OVR condition means ($R^2=0.986$, MAE $=0.017$ WER), where
\begin{equation}
    \mathrm{R}^2 = 1 - \frac{\sum_{i \in \mathrm{\{\mathrm{SNR} \times \mathrm{OVR}\}}}(\mathrm{WER}_i - \widehat{\mathrm{WER}}_i)^2}{\sum_{i \in \mathrm{\{\mathrm{SNR} \times \mathrm{OVR}\}}}(\mathrm{WER}_i - \overline{\mathrm{WER}})^2},
\end{equation}
\begin{equation}
    \mathrm{MAE} = \frac{1}{16}\sum_{i \in \mathrm{\{\mathrm{SNR} \times \mathrm{OVR}\}}}|\mathrm{WER}_i - \widehat{\mathrm{WER}}_i|,
\end{equation}
where $\mathrm{R}^2$ is the coefficient of determination, $\mathrm{MAE}$ is the mean absolute error, and $\overline{\mathrm{WER}}$ is the mean WER across the 16 conditions. 
The clip-level fit is lower ($R^2=0.433$, MAE $=0.148$ WER) because individual clips vary substantially in lexical content and acoustic difficulty. The result shows that the OVR dominates the WER increase as it has a larger coefficient than the noise term. However, the negative interaction term suggests that there is a antagonistic interaction between noise and overlap, suggesting that the effect of noise and overlap reduce each other's impact.

\subsection{Summary}
The results show a positive relationship between WER and overlap ratio, and between WER and noise (noiser conditions have higher WER than clean condition). The effect of overlap ratio is more significant in noisier conditions, while the effect of noise is more significant in lower overlap conditions. The fitted equation suggests that the effect of overlap ratio is more significant than the effect of noise, but they interact in an antagonistic way, where they reduce each other's impact on WER.


\section{RQ2: Comparison of ASR Models}
\label{sec:rq2_comparison_of_asr_models}

\begin{figure}[htbp]
    \centering
    \begin{subfigure}[t]{0.49\linewidth}
        \centering
        \includefigure[width=\linewidth]{figures/stat_cross_model_wer_vs_ovr.png}
        \caption{WER against overlap ratio, stratified by SNR.}
        \label{fig:stat_cross_model_wer_vs_ovr}
    \end{subfigure}
    \hfill
    \begin{subfigure}[t]{0.49\linewidth}
        \centering
        \includefigure[width=\linewidth]{figures/stat_cross_model_wer_vs_snr.png}
        \caption{WER against SNR condition, stratified by OVR.}
        \label{fig:stat_cross_model_wer_vs_snr}
    \end{subfigure}
    \caption{Mean WER trends across the cross-model subset. Each panel includes condition-specific lines and an overall pooled mean.}
    \label{fig:stat_cross_model_wer_trends}
\end{figure}

\subsection{Overall Comparison}
\Cref{tab:cross_model_wer_first} records the WER, ORC-WER and DSS-WER for the four ASR models on the matched 700-clip subset. The results show that \texttt{faster-whisper} has the lowest mean WER of  0.159, while \texttt{Parakeet} has the highest mean WER of 0.503, with a substantial gap between the second-highest model \texttt{wav2vec-2.0} of 0.228. The ORC-WER and DSS-WER values are generally consistent with the WER ranking, with \texttt{faster-whisper} having the lowest DSS-WER.

\Cref{fig:stat_cross_model_wer_trends} shows the WER trends against overlap ratio and SNR for the four ASR models. Both shows the same ranking of the models across conditions, with $\texttt{faster-whisper} < \texttt{WhisperX} < \texttt{wav2vec-2.0} < \texttt{Parakeet}$, which is consistent with the overall WER ranking in \Cref{tab:cross_model_wer_first}. The gap between the models remain relatively consistent as overlap ratio increases, but the gap decreases as SNR decreases. It is interesting to note that the performance of \texttt{Parakeet} is counter-intuitively better as noise increases, which might be due to difference in pretraining data and architecture. In addition, it has a dip in WER (slight increase in WER) from 0.14 OVR to 0.20 OVR. Both of these behaviours are not observed in the other three models, which might suggest that \texttt{Parakeet} and the architectural design (conformer and transducer-based architecture) is more resistant to noise. This is an interesting finding that warrants further experimentation and analysis, but it is not the main focus of this dissertation, so it will not be further investigated in this chapter.

\subsection{Extreme Cases}
This section analyses the extreme cases of the worst performing clip for each model. 

\paragraph{Worst WER for \texttt{faster-whisper}.}
Clip ID: \texttt{mix\_0004361\_0.40\_2\_7.4\_T}; WER: 1.033; Condition: OVR $=0.40$, SNR $=7.4$, noise type $=T$.

\begin{asrexample}
\footnotesize
\textbf{Hypothesis (excerpt):}
The project ... the director. Thank you for watching. If you like this video, please subscribe to the channel. ...
\end{asrexample}
In the example above, after the relevant hypothesis is finished, the \"Thank you for watching. If you like this video, please subscribe to the channel.\" are repeated twice and then \"If you like this video, please subscribe to the channel.\" is repeated 11 more times. It is worth noting that this specific repeating sentence or anything similar is not in any part of the audio. Also, every \"Thank you for watching.\" and \"If you like this video, please subscribe to the channel.\" are output as standalone segments. So this sentence is purely hallucination. Hence, the hypothesis exhibits severe hallucination and repetition, which is likely due to the LLM-based decoding strategy of \texttt{faster-whisper} that is more prone to halucination, especially under difficult conditions. If the hallucinated tokens (all prefixing "Thank you for watching. If you like this video, please subscribe to the channel.") are removed, the WER would be reduced to 0.509, which is still a high WER but much more in line with the other results in the same condition. After inspecting the worst outliers for \texttt{faster-whisper} (3 worst being \texttt{mix\_0004361\_0.40\_2\_7.4\_T}, \texttt{mix\_0003509\_0.20\_2\_7.4\_T}, \texttt{mix\_0001484\_0.14\_2\_0\_T}), it is found that they all contain severe hallucination. However, there is no clear pattern in terms of the condition of the clips, as they varied across different SNR and OVR conditions.
%to be continue

\paragraph{Worst WER for \texttt{WhisperX}.}
Clip ID: \texttt{mix\_0002729\_0.20\_2\_7.4\_T}; WER: 0.818; Condition: OVR $=0.20$, SNR $=7.4$, noise type $=T$.
\begin{asrexample}
\footnotesize
\textbf{Hypothesis (excerpt):}
A week later ... need of him.\" From that day forth, Duke was never referred to. ... From that day forth and equipped with ...
\end{asrexample}

In this example, the setence \"From that day forth, Duke was never referred to." is repeated 18 times but it came back to the correct sentence \"From that day forth and equipped with ...\" at the end. The whole hallucinated part is part of the same segment of hypothesis output by the model, which is different from the hallucination in \texttt{faster-whisper}. This suggests that the hallucination might be due to the decoding strategy of \texttt{WhisperX}, where \texttt{faster-whisper} and \texttt{WhisperX} might have different decoding strategies that lead to different pattern of hallucination. If the 